\begin{helper}[A two-color path component is not a full cycle component]
Let (z) be a node of a two-color incidence graph.  If the component through
(z) is a path component, then it is not a two-color component cycle whose
closed walk has the whole component as its support.
\end{helper}
\begin{proof}
By \noderef{TwoColorPathComponent}, the component contains a node (x)
whose neighbor set is a subsingleton.  Suppose for a contradiction that
\noderef{TwoColorComponentCycle} supplies a simple closed walk (C) whose
support is the entire component.  Then (x) lies on (C).  Rotate (C) so
that its initial and terminal node is (x).  Rotation preserves simplicity
and closedness.  The first node after (x) and the penultimate node before
the final (x) are both neighbors of (x).  They are distinct: otherwise
the rotated closed walk would repeat that node in its cyclic support, whereas
a simple cycle has no such repetition (equivalently, its first and
penultimate neighbors are distinct).  Thus the neighbor set of (x)
contains two distinct nodes, contradicting that it is a subsingleton.
\end{proof}
